\documentclass[prd,twocolumn,superscriptaddress,nofootinbib]{revtex4-2}
\usepackage{amsmath,amssymb}
\usepackage{graphicx}
\usepackage{hyperref}
\usepackage{booktabs}
\usepackage{braket}
\usepackage{bm}
\usepackage{xcolor}

% Custom macros
\newcommand{\Tr}{\mathrm{Tr}}
\newcommand{\dd}{\mathrm{d}}
\newcommand{\order}[1]{\mathcal{O}(#1)}
\newcommand{\vev}[1]{\langle #1 \rangle}

\begin{document}

\title{Sector selectivity of Koide-type mass relations from the sBootstrap}

\author{Alejandro Rivero}
\email{arivero@unizar.es}
\affiliation{BIFI, Universidad de Zaragoza, E-50018 Zaragoza, Spain}

\date{\today}

\begin{abstract}
The Koide mass relation $Q = (\sum\sqrt{m_i})^2/\sum m_i = 3/2$ holds
for charged leptons to $10^{-5}$ precision, and near-Koide tuples appear
among pseudoscalar mesons and certain quark triples, but not among vector
mesons or baryons.  This selectivity pattern has lacked a structural
explanation.  I show that the sBootstrap framework, which identifies
sfermions with known mesons and diquarks via SU(5) flavor symmetry,
provides such an explanation.  The sBootstrap mass formula from the SO(32)
embedding (conditions $z_1+z_2+z_3=0$ and $z_0^2=\sum z_k^2/3$) is proved
algebraically equivalent to the Koide condition $Q=3/2$, with the two angular
parametrizations differing by exactly $\pi/6$.  The sfermion identification
selects spin-0 two-body composites---precisely the pseudoscalar
mesons---while excluding vectors (wrong spin) and baryons (wrong composite
structure).  Through a five-link algebraic bridge---lepton Koide condition,
slepton\,=\,meson identification, shared charge structure, spin-selection
rule, and chiral/heavy-light compatibility---the framework connects the
lepton Koide relation to the hadronic sector without invoking purely
hadronic mechanisms.  It does not, however, derive the Koide angle from
first principles:\ the angle $\delta$ remains a free parameter per sector.
Nor does it derive the inverse-mass relation
$Q(1/m_d,1/m_s,1/m_b)\approx 3/2$: the turtles classify $(d,s,b)$ as
a distinguished triple, but the $Z_3$ parametrization does not close
under inversion.  The observed ratio
$\delta_{scb}\approx 3\,\delta_{e\mu\tau}$ and the discriminant coincidence
$\Delta_{\text{II}}=9=k_d^2$ remain open problems.
Three testable predictions are identified:
spin-selection (only spin-0 mesons can show Koide structure),
$N=3$ uniqueness (the Koide chain terminates at the top quark), and
the existence of charge-$\pm 4/3$ diquarks forming a Koide triple.
\end{abstract}

\maketitle

%% ===================================================================
\section{Introduction}
\label{sec:intro}
%% ===================================================================

In 1983, Koide observed that the charged lepton pole masses satisfy
\begin{equation}
  Q \;\equiv\; \frac{(\sqrt{m_e}+\sqrt{m_\mu}+\sqrt{m_\tau})^2}
  {m_e + m_\mu + m_\tau}
  \;=\; \frac{3}{2}
  \label{eq:koide}
\end{equation}
to a precision of $10^{-5}$~\cite{Koide1983prd,Koide1983plb}.  Foot
recast this as a geometric statement: the vector
$(\sqrt{m_e},\sqrt{m_\mu},\sqrt{m_\tau})$ makes an angle of exactly
$45^\circ$ with the democratic direction
$(1,1,1)$~\cite{Foot1994}.  This model-independent reformulation
separated the empirical fact from its composite-model origin and prompted
a sustained search for extensions and
mechanisms~\cite{RiveroGsponer2005}.

Rodejohann and Zhang~\cite{RodejohannZhang2011} noticed that the
mixed-charge quark triple $(c,b,t)$ also satisfies Eq.~\eqref{eq:koide}
approximately, and Cao gave the first published account~\cite{Cao2012}.
I subsequently identified the full chain
$(0,s,c)\to(-s,c,b)\to(c,b,t)$, where each link is a Koide triple with
a sign-flipped lightest charge, and showed that the chain angle satisfies
$\delta_{scb}\approx 3\,\delta_{e\mu\tau}$~\cite{Rivero2013chain}.
Zenczykowski independently confirmed this relation using the $Z_3$
parametrization~\cite{Zenczykowski2012}.  These extensions opened the
question of whether the Koide relation reflects a principle operating
across particle sectors or a collection of numerical coincidences.

A systematic scan of particle mass triples reveals a striking selectivity
pattern: pseudoscalar meson triples with a sign-flipped pion, such as
$(-\pi^0,D_s,\eta_c)$ with $|Q-3/2|=6\times 10^{-4}$, satisfy the
relation well, but no vector meson or baryon triple comes close.
The proximity of the pion and muon masses---one composite, one
elementary---was noted early~\cite{ShahMir2007} and is a foundational
clue for the supersymmetric interpretation pursued here.  The
inverse-mass tuple $K(1/m_d,1/m_s,1/m_b)$ gives $Q=1.503$,
while the inverse lepton tuple gives $Q=1.17$.  Why should
pseudoscalar mesons be selected over vectors?  Why baryons never?
Why should inverse masses of down-type quarks---and only
down-type---approach Koide?

In this paper I show that the sBootstrap~\cite{Rivero2009sboot,
Rivero2024so32} provides a structural explanation for this pattern.
The sBootstrap identifies sfermions of a supersymmetric standard model
with known mesons and diquarks via the SU(5) flavor group of five light
(``turtle'') quarks, deriving uniquely $N=3$ generations from
composite self-consistency~\cite{Rivero2007spectroscopy}, building on
the early observation by Miyazawa~\cite{Miyazawa1966} that hadronic
supersymmetry connects meson and baryon spectra, and the demonstration
by Masiero and Veneziano~\cite{MasieroVeneziano1984} that one
generation of leptons arises from SQCD with $N_c=N_f=3$.  I prove
that the sBootstrap mass formula from the SO(32)
embedding~\cite{Rivero2024so32} is algebraically equivalent to the
Koide condition $Q=3/2$, with the angle correspondence
$\alpha_{\text{Cartan}}=\delta_{Z_3}+\pi/6$.  The slepton\,=\,meson
identification then explains sector selectivity without additional
assumptions: spin-0 mesons are sfermions and inherit the Koide
constraint; spin-1 mesons and three-body baryons do not.

I am explicit about what the framework does not achieve.  The inverse
down-type relation $Q(1/m_d,1/m_s,1/m_b)\approx 3/2$ is \emph{not}
derived: the $Z_3$ parametrization does not close under mass
inversion, and the coincidence $\Delta_{\text{II}}=9=k_d^2$ connecting
the Koide discriminant to the turtle count has no known mechanism.  The
angle ratio $\delta_{scb}/\delta_{e\mu\tau}\approx 3$ is compatible
with the SU(5) structure but not derived from it.  The framework
provides a structural map, not a dynamical theory.

The paper is organized as follows.  Section~\ref{sec:sbootstrap}
reviews the sBootstrap framework and the SO(32) mass formula.
Section~\ref{sec:equivalence} presents the algebraic proof that this
mass formula implies $Q=3/2$ and derives the angle correspondence.
Section~\ref{sec:selectivity} explains the sector selectivity pattern,
including an honest assessment of the inverse down-type relation.
Section~\ref{sec:bridge} presents the five-link non-hadronic bridge
from leptons to mesons.
Section~\ref{sec:discussion} assesses what the framework explains and
what it does not.

%% ===================================================================
\section{The sBootstrap framework}
\label{sec:sbootstrap}
%% ===================================================================

\subsection{Sfermion\,=\,hadron identification}
\label{sec:identification}

The starting observation~\cite{Rivero2007spectroscopy,Rivero2009sboot}
is that the global flavor symmetry of $N_f=5$ light quarks is SU(5),
decomposed as $\text{SU}(3)_D \times \text{SU}(2)_U$ under the
separation into three down-type quarks $(d,s,b)$ and two up-type quarks
$(u,c)$.  The top quark does not form light composites and plays no role
in the flavor group; in the sBootstrap language it is the ``elephant,''
while the five participating quarks are the ``turtles.''

The two-body composites decompose under SU(5) as follows.  Meson-like
(quark--antiquark) composites transform as
\begin{equation}
  \mathbf{5}\otimes\bar{\mathbf{5}} = \mathbf{24}\oplus\mathbf{1}\,,
  \label{eq:mesons}
\end{equation}
and diquark (quark--quark) composites as
\begin{equation}
  \bar{\mathbf{5}}\otimes\bar{\mathbf{5}}
  = \mathbf{15}\oplus\mathbf{10}\,.
  \label{eq:diquarks}
\end{equation}
The identification is:
\begin{itemize}
\item \textbf{Sleptons} $\leftrightarrow$ \textbf{pseudoscalar mesons}:
  the $\mathbf{24}$ contains the scalar partners of charged leptons,
  identified with the pseudoscalar meson nonet (three generations of
  charged sleptons $\leftrightarrow$ three families of charged mesons).
\item \textbf{Squarks} $\leftrightarrow$ \textbf{diquarks}:
  the $\mathbf{15}$ and $\overline{\mathbf{15}}$ contain scalar
  partners of quarks, identified with diquark states.
\item \textbf{Sneutrinos} $\leftrightarrow$ \textbf{neutral pseudoscalars}:
  the singlet $\mathbf{1}$ accommodates the neutral sector
  ($\eta,\eta'$).
\end{itemize}
The $\mathbf{15}$ also produces three exotic diquarks of electric charge
$\pm 4/3$, one per generation, from the $(1,3)$ component of the SU(5)
decomposition~\cite{Rivero2011diquark}.

\subsection{Unique closure at $N=3$}
\label{sec:closure}

The identification above is not merely a relabeling.  The requirement
that the composite scalars reproduce exactly the sfermion content of an
$N$-generation supersymmetric standard model imposes
\begin{align}
  2N &= k_u\,k_d\,, \label{eq:closure1}\\
  2N &= \frac{k_d(k_d+1)}{2}\,, \label{eq:closure2}\\
  4N &= k_u^2 + k_d^2 - 1\,, \label{eq:closure3}
\end{align}
where $k_u$ and $k_d$ are the numbers of up-type and down-type turtle
quarks~\cite{Rivero2007spectroscopy,Rivero2024so32}.  The unique
solution is
\begin{equation}
  N=3,\qquad k_u=2,\qquad k_d=3\,.
  \label{eq:N3}
\end{equation}
This is the ``turtles all the way down'' postulate: the composites of
the turtles must be the superpartners of the turtles themselves, and
this self-consistency condition uniquely requires three generations with
five turtle quarks.

\subsection{The SO(32) embedding and mass formula}
\label{sec:so32}

The five-flavor SU(5) sits naturally inside the SO(32) group of Type~I
superstrings~\cite{Rivero2024so32}.  The adjoint of SO(32) decomposes
under $\text{SU}(5)\times\text{SU}(3)_c$ as
\begin{equation}
  496 = (\mathbf{24},\mathbf{1})
  \oplus (\mathbf{15},\bar{\mathbf{3}})
  \oplus (\overline{\mathbf{15}},\mathbf{3})
  \oplus \cdots\,,
  \label{eq:so32decomp}
\end{equation}
where the displayed terms contain exactly the scalars of a
three-generation SUSY model.

The classical energy of a pair composite with preon charges $q_a,q_b$ is
\begin{equation}
  E(q_a,q_b) = (q_a + q_b)^2\,K_\Omega\,,
  \label{eq:pair-energy}
\end{equation}
where $K_\Omega$ is a common kinematic factor.  For three generations,
the abelian charges $z_k$ ($k=1,2,3$) and a central charge $z_0$ obey
two conditions:
\begin{align}
  z_1 + z_2 + z_3 &= 0\,, \label{eq:traceless}\\
  z_0^2 &= \frac{z_1^2+z_2^2+z_3^2}{3}\,. \label{eq:normalization}
\end{align}
The masses of the three generations are then
\begin{equation}
  m_k = (z_k + z_0)^2\,.
  \label{eq:mass-formula}
\end{equation}
The SU(3) Cartan parametrization of the traceless charges is
\begin{align}
  z_1 &= \tfrac{1}{2}\cos\alpha
        + \tfrac{1}{2\sqrt{3}}\sin\alpha\,,\nonumber\\
  z_2 &= -\tfrac{1}{\sqrt{3}}\sin\alpha\,,\label{eq:cartan-param}\\
  z_3 &= -\tfrac{1}{2}\cos\alpha
        + \tfrac{1}{2\sqrt{3}}\sin\alpha\,,\nonumber
\end{align}
with $z_0 = \pm 1/\sqrt{6}$.  For charged leptons, the angle
$\alpha=0.7458$\,rad reproduces the experimental
masses~\cite{Rivero2024so32}.

%% ===================================================================
\section{Algebraic equivalence with Koide}
\label{sec:equivalence}
%% ===================================================================

\subsection{Proof that $Q=3/2$}
\label{sec:proof}

I now show that the conditions~\eqref{eq:traceless}
and~\eqref{eq:normalization} imply the Koide relation $Q=3/2$ for the
masses~\eqref{eq:mass-formula}.

\medskip\noindent\textit{Numerator.}\quad
Taking square roots of the masses (all positive by construction):
\begin{equation}
  \sqrt{m_k} = z_k + z_0\,,
\end{equation}
so
\begin{equation}
  \sum_{k=1}^3 \sqrt{m_k}
  = (z_1+z_2+z_3) + 3z_0
  = 0 + 3z_0
  = 3z_0\,,
  \label{eq:numerator-step}
\end{equation}
using the tracelessness condition~\eqref{eq:traceless}.  Hence
\begin{equation}
  \Bigl(\sum_k\sqrt{m_k}\Bigr)^2 = 9\,z_0^2\,.
  \label{eq:numerator}
\end{equation}

\medskip\noindent\textit{Denominator.}\quad
Expanding the masses:
\begin{align}
  \sum_{k=1}^3 m_k
  &= \sum_k(z_k+z_0)^2\nonumber\\
  &= \sum_k z_k^2 + 2z_0\sum_k z_k + 3z_0^2\nonumber\\
  &= \sum_k z_k^2 + 0 + 3z_0^2\,.
  \label{eq:denom-step1}
\end{align}
Applying the normalization condition~\eqref{eq:normalization},
$\sum_k z_k^2 = 3z_0^2$, yields
\begin{equation}
  \sum_k m_k = 3z_0^2 + 3z_0^2 = 6z_0^2\,.
  \label{eq:denominator}
\end{equation}

\medskip\noindent\textit{Ratio.}\quad
Dividing~\eqref{eq:numerator} by~\eqref{eq:denominator}:
\begin{equation}
  \boxed{Q = \frac{9z_0^2}{6z_0^2} = \frac{3}{2}\,.}
  \label{eq:Q-proof}
\end{equation}
This holds for any value of $z_0\neq 0$ and any angle in
the parametrization~\eqref{eq:cartan-param}.  The Koide relation is
not an additional constraint imposed on the sBootstrap mass
formula---it is an algebraic \emph{consequence} of the tracelessness
and normalization conditions, Eqs.~\eqref{eq:traceless}
and~\eqref{eq:normalization}.

\subsection{Connection to Koide's original condition}
\label{sec:koide-condition}

In Koide's 1981 model, the charged lepton mass matrix is written as
$M_l=(U+V)^2$, where $U\propto\mathbf{1}$ is the central part and $V$
is traceless, with the constraint $\Tr[U^2]=\Tr[V^2]$~\cite{Koide1983prd}.
In the sBootstrap parametrization, $U=z_0\,\mathbf{1}$ and
$V=\text{diag}(z_1,z_2,z_3)$, so
\begin{equation}
  \Tr[U^2] = 3z_0^2\,,\qquad
  \Tr[V^2] = z_1^2+z_2^2+z_3^2\,.
  \label{eq:traces}
\end{equation}
The normalization condition~\eqref{eq:normalization} is therefore
identical to Koide's constraint:
\begin{equation}
  z_0^2 = \frac{\sum_k z_k^2}{3}
  \quad\Longleftrightarrow\quad
  \Tr[U^2] = \Tr[V^2]\,.
  \label{eq:koide-sboot-equiv}
\end{equation}
The sBootstrap does not introduce a new mass formula.  It provides a
physical context---composite scalar normalization within a SUSY
framework---for the same algebraic condition that Koide identified in
1981.

\subsection{Angle correspondence}
\label{sec:angle}

The Koide chain literature uses the $Z_3$ phase parametrization
\begin{equation}
  q_k = \sqrt{M_0}\left(1+\sqrt{2}\,\cos\!\left(\frac{2\pi k}{3}
  +\delta\right)\right),\quad k=1,2,3\,,
  \label{eq:z3-param}
\end{equation}
with masses $m_k=q_k^2$ and $\sqrt{M_0}$ setting the overall
scale~\cite{Brannen2006,Zenczykowski2012}.  Comparing with
the sBootstrap masses $m_k=(z_k+z_0)^2$, the identification is
\begin{equation}
  z_0 = \sqrt{M_0}\,,\qquad
  z_k = \sqrt{2\,M_0}\,\cos\!\left(\frac{2\pi k}{3}+\delta\right).
  \label{eq:z-q-map}
\end{equation}
These $z_k$ satisfy~\eqref{eq:traceless} identically (the sum of three
equally spaced cosines vanishes), and~\eqref{eq:normalization} by
construction (the sum of their squares equals $3M_0/2\cdot 2 = 3M_0$,
while $3z_0^2=3M_0$, recalling the normalization $z_0=\pm 1/\sqrt{6}$
in the unit-norm convention vs.\ $z_0=\sqrt{M_0}$ in physical units).

The coordinate transformation between the SU(3) Cartan
basis~\eqref{eq:cartan-param} and the $Z_3$ phase
basis~\eqref{eq:z3-param} is a rotation by $\pi/6$:
\begin{equation}
  \alpha_{\text{Cartan}} = \delta_{Z_3} + \frac{\pi}{6}\,.
  \label{eq:angle-shift}
\end{equation}
This can be verified by writing out both parametrizations
and matching the three charge components, using
$\cos(\theta+\pi/6)=(\sqrt{3}\cos\theta-\sin\theta)/2$.

\medskip\noindent\textit{Numerical check.}\quad For charged leptons:
\begin{alignat}{2}
  &\text{SO(32) paper:} &\quad \alpha &= 0.7458\;\text{rad}
  = 42.73^\circ\,,\nonumber\\
  &\text{Koide chain:}  &\quad \delta &= 0.2222\;\text{rad}
  = 12.73^\circ\,,\label{eq:angle-check}\\
  &\text{Difference:}   &\quad \alpha-\delta &= 30.00^\circ
  = \frac{\pi}{6}\;.\nonumber
\end{alignat}
The correspondence is exact.  The two parametrizations describe the same
one-parameter family of mass spectra, rotated by $\pi/6$ in the charge
space on the SU(3) Cartan torus.  The value $\delta=2/9$\,rad was first
identified in Brannen's circulant mass matrix
construction~\cite{Brannen2006}, which is mathematically identical to
the $Z_3$ parametrization.

%% ===================================================================
\section{Sector selectivity from the sBootstrap}
\label{sec:selectivity}
%% ===================================================================

The sBootstrap identification of sfermions with specific composites
provides a structural mechanism for the observed selectivity of
Koide-type relations.  I analyze each sector in turn.

\subsection{Charged leptons}
\label{sec:leptons}

Charged leptons are the fundamental fermions of the model.  Their
scalar superpartners (sleptons) are identified with the charged
pseudoscalar mesons living in the $\mathbf{24}$ of SU(5).  The Koide
condition $\Tr[U^2]=\Tr[V^2]$, Eq.~\eqref{eq:koide-sboot-equiv}, is
interpreted as a SUSY condition on the (lepton, slepton\,=\,meson)
multiplet: it constrains the mass matrix of both the fermion and its
scalar partner simultaneously.  That the charged lepton masses satisfy
$Q=3/2$ to $10^{-5}$ precision is, in this picture, a consequence of the
SUSY condition being (approximately) preserved.

Sumino~\cite{Sumino2009} has shown how a $\text{U}(3)\times\text{O}(3)$
family gauge symmetry can protect the Koide relation against QED
radiative corrections at the pole-mass level.  The sBootstrap framework
is compatible with this mechanism: the family gauge symmetry would
operate within the lepton sector, while the sBootstrap provides the
structural context linking leptons to mesons.

\subsection{Pseudoscalar mesons}
\label{sec:pseudoscalars}

Pseudoscalar mesons ($J^P=0^-$) are the lightest spin-0 $q\bar{q}$
composites.  In the sBootstrap, they \emph{are} the sleptons: the
$\mathbf{24}$ of SU(5) from the $\mathbf{5}\otimes\bar{\mathbf{5}}$
decomposition.  The Koide condition on sleptons translates directly to a
constraint on pseudoscalar meson mass triples.

The numerical evidence confirms this: several pseudoscalar meson triples
satisfy $Q\approx 3/2$ with a sign-flipped pion~\cite{Rivero2007spectroscopy},
as shown in Table~\ref{tab:sectors}.
The best is $(-\pi^0,D_s,\eta_c)$ with $|Q-3/2|=6\times 10^{-4}$.
The sBootstrap explains \emph{why} these tuples exist:
pseudoscalar mesons show near-Koide behavior because they are the
sfermions, subject to the same SUSY-origin constraint as the leptons.

The match is ``near'' rather than exact because
SUSY is not an exact symmetry at hadronic scales: QCD dynamics, chiral
symmetry breaking, and quark mass differences all contribute to
deviations from the $\Tr[U^2]=\Tr[V^2]$ condition.  The sBootstrap
predicts that the deviations should be controlled by the ratio of
QCD effects to the slepton mass scale.

\subsection{Vector mesons}
\label{sec:vectors}

Vector mesons ($J^P=1^-$) are spin-1 $q\bar{q}$ states.  They are
\emph{not} sfermions.  In any supersymmetric framework, the scalar
partners of fermions carry spin~0.  The sBootstrap SUSY operates
exclusively on scalar composites: only spin-0 mesons (pseudoscalar and,
potentially, scalar) qualify as sfermions.  No Koide-type constraint
applies to vectors.

The data confirm this: the best vector meson triple achieves only
$Q\approx 1.31$, far from $3/2$.

\subsection{Baryons}
\label{sec:baryons}

Baryons are $qqq$ (three-body) composites.  The sBootstrap SUSY
operates on $q\bar{q}$ (meson) and $qq$ (diquark) composites---both
two-body.  Three-body states fall outside the
$\mathbf{5}\otimes\bar{\mathbf{5}}$ and
$\bar{\mathbf{5}}\otimes\bar{\mathbf{5}}$ decompositions.  No Koide
constraint applies.

Numerically, the best baryon triple gives $Q\approx 1.19$.

\subsection{The quark chain}
\label{sec:chain}

The five turtle quarks $(u,c,d,s,b)$ are the fundamental objects of the
SU(5) flavor group.  The Koide chain
$(0,s,c)\to(-s,c,b)\to(c,b,t)$~\cite{Rivero2013chain} connects
successive quark triples, with each link satisfying $Q\approx 3/2$.
The SU(5) structure is \emph{compatible} with this chain: the five
turtles span exactly the quarks appearing in the chain, and the top
(the elephant) serves as the chain endpoint.

However, the factor-of-three angle relation
$\delta_{scb}\approx 3\,\delta_{e\mu\tau}$~\cite{Zenczykowski2012,
Rivero2013chain} is not derived from the SU(5) turtle structure alone.
The sBootstrap explains \emph{which} quarks participate (the five
turtles) but not \emph{why} the chain angles are related to the lepton
angle by a factor of three.

\subsection{Inverse down-type quarks: an honest negative}
\label{sec:inverse}

The tuple $K(1/m_d,1/m_s,1/m_b)$ gives $Q=1.503$, while
$K(1/m_e,1/m_\mu,1/m_\tau)$ gives $Q=1.17$ and
$K(1/m_u,1/m_c,1/m_t)$ gives $Q=1.09$.  Only the down-type inverse
approaches Koide.

The sBootstrap provides a \emph{structural reason} for why $(d,s,b)$ is
distinguished among all quark triples: the closure
condition~\eqref{eq:closure2} requires exactly $k_d=3$ down-type
turtles, and the Standard Model contains exactly three down-type quarks.
The bootstrap equations do not select this triple by fitting it to the
data; they select it algebraically through the self-consistency of the
composite spectrum.

However, the inverse-mass Koide relation itself is \emph{not} derived
from the sBootstrap.  Three distinct arguments show this:

\begin{enumerate}
\item \textbf{The $Z_3$ parametrization does not close under inversion.}
The inverse charges $w_k = 1/q_k = 1/[\sqrt{M_0}(1+\sqrt{2}\cos(2\pi
k/3+\delta))]$ are not of the $Z_3$ form $A(1+\sqrt{2}\cos(2\pi
k/3+\delta'))$ for any constants $A$ and $\delta'$.  The Koide condition
$Q=3/2$ is automatic for direct masses in the $Z_3$ parametrization but
\emph{not} for inverse masses.

\item \textbf{No Cartan algebra involution exists} that maps $m_k\to
1/m_k$ within the SO(32) mass formula.  The masses $m_k=(z_k+z_0)^2$
are quadratic in the charges; the inverse $1/m_k = 1/(z_k+z_0)^2$ is
not related to any linear transformation of the Cartan components.

\item \textbf{The diquark channel provides no constraint.}  Down-type
diquarks in the $(\mathbf{6},\mathbf{1})$ component of $\mathbf{15}$
have masses $E(q_a,q_b)=(q_a+q_b)^2 K_\Omega$ that depend on
\emph{sums} of charges.  These sum constraints do not impose conditions
on \emph{individual} inverse masses $1/m_k$.
\end{enumerate}

\noindent There is a curious numerical coincidence.  The discriminant of
the inverse-Koide quadratic at $\delta=15^\circ$ evaluates to
$\Delta_{\text{II}}=9$, and $\sqrt{\Delta_{\text{II}}}=3=k_d$.  Whether
this connection between the Koide discriminant and the turtle count is
accidental or points to a deeper structure is an open question.  I note
that no representation-theoretic quantity of SU(5) (dimensions, Casimir
invariants, Dynkin indices of the $\mathbf{24}$ or $\mathbf{15}$)
produces the factor 3 in a natural combination; the only quantity that
equals~3 is $k_d$ itself.

The status of the inverse down-type tuple is therefore: \emph{structurally
selected triple, unexplained mass relation}.

\subsection{Sector comparison}
\label{sec:comparison}

Table~\ref{tab:sectors} summarizes the sector selectivity analysis.  The
sBootstrap provides a complete structural explanation for four of six
sectors and partial structural compatibility for the remaining two.
There are zero contradictions with the empirical data.

\begin{table*}[t]
\centering
\caption{Sector selectivity: empirical Koide status and sBootstrap
  explanation.  The Koide ratio $Q$ is computed from PDG~2024 masses
  in the $m_q(m_q)$ convention for quarks and pole masses for
  leptons~\cite{PDG2024}.  The sBootstrap prediction column indicates
  whether the framework selects or excludes each sector and why.}
\label{tab:sectors}
\begin{tabular}{@{}llcccl@{}}
\toprule
Sector & Best tuple & $Q$ & $|Q-3/2|$
  & Selected? & sBootstrap reason \\
\midrule
Charged leptons & $(e,\mu,\tau)$
  & $1.50000$ & $2\times 10^{-5}$
  & Yes & SUSY condition on (lepton, meson) multiplet \\
Pseudoscalar mesons & $(-\pi^0,D_s,\eta_c)$
  & $1.5006$ & $6\times 10^{-4}$
  & Yes & Mesons \emph{are} the sleptons (spin-0, $q\bar{q}$) \\
Quark chain & $(0,s,c)$
  & $1.505$ & $5\times 10^{-3}$
  & Partial & Compatible with SU(5) turtles; angle ratio not derived \\
Inverse down-type & $K(1/d,1/s,1/b)$
  & $1.503$ & $3\times 10^{-3}$
  & Partial & $(d,s,b)$ selected as SU(3)$_D$; inversion not derived \\
Vector mesons & $(-\rho,B_s^*,\Upsilon)$
  & $1.31$ & $0.19$
  & No & Wrong spin ($J=1$, not sfermions) \\
Baryons & $(-p,\Sigma_b^-,\Omega_b^-)$
  & $1.19$ & $0.31$
  & No & Wrong composite structure (3-body, not 2-body) \\
\bottomrule
\end{tabular}
\end{table*}

%% ===================================================================
\section{The non-hadronic bridge}
\label{sec:bridge}
%% ===================================================================

The central structural contribution of the sBootstrap to the
Koide problem is a five-link algebraic bridge connecting two apparently
unrelated observations: the charged lepton Koide relation ($Q=3/2$
exact) and the near-Koide behavior of pseudoscalar meson triples.  I
present each link explicitly, with its algebraic content and confidence
assessment.

\subsection{The five-link bridge}
\label{sec:five-links}

\begin{enumerate}
\item \textbf{Link~1: Lepton Koide = trace constraint.}
The lepton mass matrix $M_l=(U+V)^2$ with $\Tr[U^2]=\Tr[V^2]$ is
algebraically identical to the sBootstrap conditions
$z_1+z_2+z_3=0$, $z_0^2=\sum z_k^2/3$.  This constrains the mass
spectrum to a one-parameter $Z_3$ orbit, with the Koide condition
$Q=3/2$ as an automatic consequence
(Sec.~\ref{sec:proof}).

\item \textbf{Link~2: Sleptons = mesons in the $\mathbf{24}$ of
SU(5).}
With five turtles $\{u,c,d,s,b\}$ and one elephant $\{t\}$, the
SU(5) flavor products
$\mathbf{5}\otimes\bar{\mathbf{5}}=\mathbf{24}\oplus\mathbf{1}$
identify $q\bar{q}$ mesons with sleptons.  This is a
representation-theoretic identification, not a
fit~\cite{Rivero2009sboot}.

\item \textbf{Link~3: Shared charge structure transfers the
constraint.}
The \emph{same} set of charges $\{z_k\}$ from the SU(3) Cartan
subalgebra determines both lepton masses
$m_k=(z_k+z_0)^2$ and meson pair energies
$E(q_a,q_b)=(q_a+q_b)^2 K_\Omega$~\cite{Rivero2024so32}.  The Koide
condition constraining the lepton sector transfers to the meson sector
through these shared charges.  The transfer is not exact because QCD
corrections modify the meson masses, producing the observed
$|Q-3/2|\sim 10^{-4}\text{--}10^{-3}$ deviations.

\item \textbf{Link~4: Spin-selection and bilinear exclusion.}
Sfermions are spin-0 scalars.  This selects pseudoscalar mesons
($J^P=0^-$, the lightest spin-0 $q\bar{q}$ states) and excludes
vector mesons ($J^P=1^-$, spin-1) by an algebraic selection rule.
Baryons ($qqq$) are trilinear composites that do not appear in the
bilinear $\mathbf{5}\otimes\bar{\mathbf{5}}$ products and are
excluded on separate grounds.

\item \textbf{Link~5: Compatibility with the chiral/heavy-light
selector.}
The sBootstrap's spin-0 selection provides a structural explanation for
the chiral and heavy-light selectors identified in purely hadronic
analyses: chiral symmetry breaking affects pseudoscalars (which are the
Goldstone sector) differently from vectors, and heavier turtle quarks
$(c,b)$ have larger charges that place their meson composites closer to
the SUSY mass formula regime.  The pion sign flip arises from chiral
symmetry breaking and is a QCD correction, not an sBootstrap prediction.
\end{enumerate}

\noindent The bridge is non-hadronic: it connects leptons to mesons
through supersymmetry, not through QCD.  Previous attempts to connect
the lepton Koide relation to hadronic physics relied on purely hadronic
mechanisms (chiral symmetry, heavy-quark symmetry, confinement-scale
dynamics).  These can explain why certain meson triples approach Koide
within the hadronic sector, but they cannot explain the connection
\emph{to} the lepton relation, since leptons are not hadrons.

The hierarchy of precision---leptons closer to exact Koide than
mesons---is naturally explained: Sumino's family gauge
symmetry~\cite{Sumino2009} can protect the lepton side of the bridge
against QED radiative corrections, while the meson side receives QCD
corrections of order $\Lambda_{\text{QCD}}/m_{\text{meson}}$ that are
not cancelled.

\subsection{What the bridge does not explain}
\label{sec:gaps}

The sBootstrap bridge has clear limitations:
\begin{enumerate}
\item It does not derive the Koide angle $\delta$.  For each sector,
  $\delta$ remains a free parameter determined by the mass data.  For
  leptons, $\delta_{e\mu\tau}=12.73^\circ$; for the quark chain,
  $\delta_{scb}\approx 3\times 12.73^\circ\approx 38^\circ$.
\item It does not derive the ratio
  $\delta_{scb}/\delta_{e\mu\tau}\approx 3$, which would require
  understanding the $\text{SU}(5)\to\text{SU}(3)\times\text{SU}(2)$
  breaking pattern at a dynamical level.
\item It does not derive the inverse down-type relation
  $Q(1/m_d,1/m_s,1/m_b)\approx 3/2$: the $Z_3$ parametrization does
  not close under inversion (Sec.~\ref{sec:inverse}).
\item It does not explain \emph{which} specific meson triples satisfy
  $Q\approx 3/2$: the selection of tuples like
  $(-\pi^0,D_s,\eta_c)$ requires the mass formula with a
  specific angle, which the sBootstrap does not fix.
\item It does not predict specific mass values.  The overall scale $M_0$
  and the angle $\delta$ are both determined from data, not from first
  principles.
\end{enumerate}

%% ===================================================================
\section{Discussion}
\label{sec:discussion}
%% ===================================================================

\subsection{Assessment: partial structural insight}

The milestone verdict for the sBootstrap contribution to the Koide
program is \emph{partial structural insight with a genuine non-hadronic
bridge}.  More precisely, the sBootstrap provides four things that no
other framework offers simultaneously:
\begin{enumerate}
\item A genuine non-hadronic bridge from leptons to mesons, via the
  slepton\,=\,meson identification in the $\mathbf{24}$ of SU(5).
\item A structural explanation for sector selectivity: spin-0 two-body
  composites are selected, vectors and baryons excluded.
\item An algebraic identification of the sBootstrap mass formula with
  the Koide condition, showing they are not merely compatible but
  identical (Sec.~\ref{sec:proof}).
\item A generation-counting argument: $N=3$ uniquely from the composite
  self-consistency
  conditions~\eqref{eq:closure1}--\eqref{eq:closure3}.
\end{enumerate}
These results constitute a structural insight, not a dynamical
derivation.  The sBootstrap explains \emph{which} sectors should show
Koide behavior (and why), but does not explain \emph{how} the specific
mass values arise.  The Koide angle $\delta$ encodes the mass spectrum
within each sector, and the sBootstrap leaves it as a free parameter.

\subsection{Consistency with prior exclusions}

An earlier systematic analysis of the Koide landscape established five
exclusion claims: no universal Koide mechanism; vectors fail; baryons
fail; Yukawaon/Sumino falsified as universal mechanism; hadronic-only
embedding survives as the best framework~\cite{Rivero2013chain}.  The
sBootstrap is consistent with all five.  It does not claim universal
Koide (it claims \emph{selective} Koide).  It predicts vector and baryon
failure from spin-selection and composite structure.  It is independent
of the Yukawaon mechanism but compatible with Sumino's gauge protection
as an IR mechanism.  And it extends the hadronic-only embedding by
providing the non-hadronic bridge that was previously absent.

\subsection{Community context}

A systematic evaluation of 43 ideas from extended Physics Forums
discussions on Koide-type relations~\cite{Rivero2007spectroscopy,
Brannen2006} found that 70\% are either incorporated in the sBootstrap
framework or independently support it, with zero contradictions.  Three
findings deserve special mention:

\begin{itemize}
\item Brannen's circulant mass matrix with $\delta=2/9$~\cite{Brannen2006}
  is mathematically identical to the $Z_3$/Cartan parametrization used
  here, derived independently from Clifford algebra preons.  The two
  parametrizations differ by the $\pi/6$ rotation of
  Eq.~\eqref{eq:angle-shift}.
\item The pion-mass ratio identity
  $m_{\pi^+}/m_{\pi^0} = 1 + \sqrt{m_\mu/m_Z}$
  holds to 7\,ppm with PDG~2024 masses.  This relation connects a
  hadronic mass splitting to an electroweak/lepton quantity and is
  consistent with the sBootstrap's pion-as-slepton identification, but
  is not derived from it.
\item Masiero and Veneziano's demonstration~\cite{MasieroVeneziano1984}
  that one generation of leptons arises from SQCD with $N_c=N_f=3$
  provides the strongest published support for the sBootstrap's
  mesino\,=\,lepton identification, predating the sBootstrap by two
  decades.
\end{itemize}

\subsection{Honest negatives}

Three important results are \emph{not} achieved by the sBootstrap:

\begin{enumerate}
\item \textbf{The inverse down-type relation is not derived.}  The
  turtles classify $(d,s,b)$ as a distinguished triple, but the $Z_3$
  parametrization does not close under mass inversion
  (Sec.~\ref{sec:inverse}).  The factor-of-three coincidence
  $\sqrt{\Delta_{\text{II}}}=k_d$ remains unexplained.
\item \textbf{The $3\times$ angle ratio is not derived.}  No ratio of
  SU(5) representation invariants (dimensions, Casimirs, Dynkin indices
  of $\mathbf{24}$ vs.\ $\mathbf{15}$) produces the factor~3.  The
  only quantity equal to~3 is $k_d$ itself.
\item \textbf{No specific mass values are predicted.}  The angle
  $\delta$ is free per sector, so the framework cannot predict any
  individual mass.  Given any three masses satisfying $Q=3/2$, the
  framework can accommodate them by choosing an appropriate~$\delta$.
\end{enumerate}

\subsection{Predictions}
\label{sec:predictions}

The sBootstrap makes three testable predictions regarding
Koide-type relations:
\begin{enumerate}
\item \textbf{Spin-selection.}  The sBootstrap predicts that
  Koide-type mass structure appears \emph{only} in spin-0 meson sectors
  (pseudoscalar and potentially scalar mesons, such as $f_0$, $a_0$,
  $K_0^*$) and \emph{not} in spin-1 (vectors) or three-body composites
  (baryons).  The sfermion identification is the mechanism: only spin-0
  $q\bar{q}$ composites qualify as scalar partners.  This prediction is
  structural and does not depend on the free parameter $\delta$.  Any
  vector meson triple with $Q\approx 3/2$ would falsify it.
\item \textbf{$N=3$ uniqueness and chain termination.}  The closure
  conditions~\eqref{eq:closure1}--\eqref{eq:closure3} have no solution
  for $N=4$.  The Koide chain terminates at the top quark (the
  elephant): there is no fourth-generation link.  Any discovery of a
  fourth generation of quarks, or a Koide chain continuing beyond the
  top, would falsify the sBootstrap's generation-counting argument.
\item \textbf{Exotic $\pm 4/3$ diquarks.}  The $(1,3)$ component of the
  $\mathbf{15}$ produces three exotic diquarks of charge $\pm 4/3$
  (the $uu$, $uc$, $cc$ states).  If observed, they would form their
  own Koide-type triple.  This is a conditional prediction:\ naive
  contact-interaction estimates give $Q=2.0\text{--}2.8$ for these
  states, far from $3/2$, but the diquark mass dynamics in the
  $\mathbf{15}$ may differ from the $\mathbf{24}$ where Koide is
  established.
\end{enumerate}

The structural predictions (spin-selection and $N=3$) are strongly
falsifiable: they make definite yes/no claims that do not depend on
$\delta$.  The diquark prediction is conditional on the existence of
states not yet observed.  A deferred prediction---the neutrino Koide
relation, requiring the right-handed neutrinos that the sBootstrap
demands for internal consistency---becomes testable when absolute
neutrino masses are measured.

\subsection{Comparison with other approaches}
\label{sec:other}

\textit{Yukawaon models.}---Koide's yukawaon
framework~\cite{Koide2009yukawaon} uses $\text{U}(3)\times S_3$ family
symmetry to derive mass matrices whose vacuum expectation values
reproduce the Koide relation.  This is a more dynamical approach than
the sBootstrap, but it applies sector by sector and does not address the
pseudoscalar/vector selectivity.  The sBootstrap and yukawaon approaches
are not mutually exclusive; the yukawaon could provide the dynamical
mechanism that the sBootstrap's structural framework leaves unspecified.

\textit{Superconformal light-front holographic QCD.}---The
Brodsky-de~Teramond-Dosch program~\cite{Brodsky2015superconformal}
establishes a hadronic supersymmetry connecting meson and baryon Regge
trajectories through a superconformal algebra.  This framework naturally
distinguishes pseudoscalar from vector mesons and could potentially
explain the meson Koide selectivity at a QCD level.  However, it
has not been directly connected to Koide-type sum rules, and it does not
provide a bridge to the lepton sector.  A synthesis of the sBootstrap
(lepton--meson bridge) with LFHQCD (hadronic mass patterns) is an
interesting direction for future work.

\textit{$Z_3$ parametrization.}---Zenczykowski's systematic use of the
$Z_3$ angle parametrization~\cite{Zenczykowski2012,Zenczykowski2013,
Zenczykowski2018} provided the algebraic foundation on which the chain
structure was recognized.  The sBootstrap adds the physical
interpretation of the Cartan angle
(Eq.~\eqref{eq:angle-shift}) and the generation-counting context, but
the phenomenological content of the $Z_3$ parametrization is independent
of the sBootstrap and predates it.  The discrete $S_3$ symmetry
underlying the parametrization was already implicit in the
Harari-Haut-Weyers quark mass ansatz~\cite{HarariHautWeyers1978}.

\subsection{Open questions}

Several questions remain open for future work:
\begin{itemize}
\item What dynamical mechanism fixes the Koide angle $\delta$?  The
  angle determines the mass spectrum within each sector and is the key
  quantity not derived by the sBootstrap.  Possibilities include
  the Sumino gauge sector~\cite{Sumino2009} or QCD vacuum dynamics.
\item What explains the chain ratio $\delta_{scb}\approx
  3\,\delta_{e\mu\tau}$?  The factor of three might emerge from the
  $\text{SU}(5)\to\text{SU}(3)\times\text{SU}(2)$ breaking pattern,
  but no combination of Casimir invariants or representation dimensions
  produces it.  Brannen's $\alpha$-correction formula
  $\delta=2/9-4\pi\alpha'/3^{12}$, which matches the lepton angle
  to 10~digits~\cite{Brannen2006}, suggests a perturbative origin
  for the angle that remains to be derived.
\item Can the inverse down-type relation be derived from a dual structure
  in the diquark channel?  The
  $\bar{\mathbf{5}}\otimes\bar{\mathbf{5}}=\mathbf{15}\oplus\mathbf{10}$
  decomposition contains the diquark sector; the present analysis shows
  that sum constraints on diquark masses do not constrain individual
  inverse masses (Sec.~\ref{sec:inverse}).
\item How do Yukawa corrections to quark mass running affect the
  Koide angles?  At one-loop QCD, mass ratios (and hence Koide angles)
  are exactly RG-invariant~\cite{EspositoSantorelli1995,XingZhang2006}.
  Yukawa contributions, dominated by the top, are flavor-dependent and
  could shift $\delta_{scb}$ and $\delta_w$ at the percent level.
\end{itemize}

%% ===================================================================
\section{Conclusions}
\label{sec:conclusions}
%% ===================================================================

The sBootstrap provides a partial structural explanation for the observed
sector selectivity of Koide-type mass relations---partial because the
Koide angle $\delta$ remains free, but structural because the selectivity
pattern is derived from algebraic and representation-theoretic arguments
rather than fitted to data.

Through the sfermion\,=\,hadron identification via SU(5) flavor
symmetry, the framework selects spin-0 two-body composites (pseudoscalar
mesons) as the hadronic sector where Koide-type behavior is expected,
while excluding vectors (wrong spin) and baryons (wrong composite
structure).  The algebraic proof that the sBootstrap mass formula implies
$Q=3/2$, Eqs.~\eqref{eq:traceless}--\eqref{eq:Q-proof}, establishes
that the two frameworks are not merely compatible but identical, with the
angle correspondence $\alpha_{\text{Cartan}}=\delta_{Z_3}+\pi/6$.

The five-link non-hadronic bridge---trace constraint, slepton\,=\,meson
identification, shared charge transfer, spin-selection, and
chiral/heavy-light compatibility---connects the lepton Koide relation to
hadronic mass patterns through SUSY rather than QCD.  The framework is
consistent with all six particle sectors examined: four explained
(leptons, pseudoscalar mesons, vectors, baryons), two compatible
(quark chain, inverse down-type), zero contradictions.  It makes three
testable predictions: spin-selection, $N=3$ uniqueness, and the
$\pm 4/3$ diquark sector.

What the sBootstrap does not provide is equally important.  The Koide
angle $\delta$ remains a free parameter, so no individual mass is
predicted.  The chain ratio
$\delta_{scb}/\delta_{e\mu\tau}\approx 3$ is not derived; no ratio of
SU(5) invariants produces it.  The inverse down-type relation is
structurally motivated (the turtles select $(d,s,b)$) but not proven:
the $Z_3$ parametrization does not close under mass inversion, and the
coincidence $\Delta_{\text{II}}=9=k_d^2$ has no known mechanism.  The
framework offers a structural map, not a dynamical theory.  The
derivation of $\delta$ from first principles remains the central open
problem in the Koide program.

\begin{acknowledgments}
I thank the participants of the Physics Forums community for
discussions on Koide-type relations, especially C.~Brannen for the
circulant parametrization and H.~de~Velde for the anomalous moment
coincidences.  I thank M.~Porter for identifying the Masiero-Veneziano
and Seiberg duality connections.  This work was supported by the
BIFI research infrastructure at the Universidad de Zaragoza.
\end{acknowledgments}

%% ===================================================================
%% BIBLIOGRAPHY
%% ===================================================================

\begin{thebibliography}{30}

\bibitem{Koide1983prd}
Y.~Koide,
``New view of quark and lepton mass hierarchy,''
Phys.\ Rev.\ D \textbf{28}, 252 (1983).

\bibitem{Koide1983plb}
Y.~Koide,
``A fermion-boson composite model of quarks and leptons,''
Phys.\ Lett.\ B \textbf{120}, 161 (1983).

\bibitem{Foot1994}
R.~Foot,
``A note on Koide's lepton mass relation,''
arXiv:hep-ph/9402242 (1994).

\bibitem{RodejohannZhang2011}
W.~Rodejohann and H.~Zhang,
``Extended empirical fermion mass relation,''
Phys.\ Lett.\ B \textbf{698}, 152 (2011),
arXiv:1101.5525 [hep-ph].

\bibitem{Cao2012}
F.~G.~Cao,
``Neutrino masses from lepton and quark mass relations,''
Phys.\ Rev.\ D \textbf{85}, 113003 (2012).

\bibitem{Rivero2013chain}
A.~Rivero,
``Koide equations for quark mass triplets,''
arXiv:1111.7232 [hep-ph] (2013).

\bibitem{Zenczykowski2012}
P.~Zenczykowski,
``Remark on Koide's $Z_3$-symmetric parametrization of quark masses,''
Phys.\ Rev.\ D \textbf{86}, 117303 (2012),
arXiv:1210.4125 [hep-ph].

\bibitem{Zenczykowski2013}
P.~Zenczykowski,
``Koide's $Z_3$-symmetric parametrization, quark masses and mixings,''
Phys.\ Rev.\ D \textbf{87}, 077302 (2013),
arXiv:1301.4143 [hep-ph].

\bibitem{Zenczykowski2018}
P.~Zenczykowski,
``A note on Koide's doubly special parametrization of quark masses,''
Open Physics \textbf{16}, 427 (2018).

\bibitem{Rivero2007spectroscopy}
A.~Rivero,
``Third spectroscopy with a hint of superstrings,''
arXiv:0710.1526 [hep-ph] (2007).

\bibitem{Rivero2009sboot}
A.~Rivero,
``Unbroken supersymmetry without new particles,''
arXiv:0910.4793 [hep-ph] (2009).

\bibitem{Rivero2024so32}
A.~Rivero,
``An interpretation of scalars in SO(32),''
Eur.\ Phys.\ J.\ C \textbf{84}, 1058 (2024),
arXiv:2407.05397 [hep-ph].

\bibitem{Rivero2011diquark}
A.~Rivero,
``A possible origin of the $q=4/3$ diquark,''
arXiv:1111.7230 [hep-ph] (2011).

\bibitem{Sumino2009}
Y.~Sumino,
``Family gauge symmetry as an origin of Koide's mass formula and
charged lepton spectrum,''
JHEP \textbf{0905}, 075 (2009),
arXiv:0812.2103 [hep-ph].
See also: Phys.\ Lett.\ B \textbf{671}, 477 (2009).

\bibitem{Brannen2006}
C.~Brannen,
``Koide mass equations for hadrons'' (2006, unpublished).

\bibitem{Koide2009yukawaon}
Y.~Koide,
``Charged lepton mass spectrum and supersymmetric yukawaon model,''
arXiv:0906.3370 (2009).

\bibitem{Brodsky2015superconformal}
S.~J.~Brodsky, G.~F.~de~Teramond, and H.~G.~Dosch,
``Superconformal baryon-meson symmetry and light-front holographic QCD,''
Phys.\ Rev.\ D \textbf{91}, 085016 (2015),
arXiv:1501.00959 [hep-ph].

\bibitem{HarariHautWeyers1978}
H.~Harari, H.~Haut, and J.~Weyers,
``Quark masses and Cabibbo angles,''
Phys.\ Lett.\ B \textbf{78}, 459 (1978).

\bibitem{EspositoSantorelli1995}
S.~Esposito and P.~Santorelli,
``A geometric picture for fermion masses,''
Mod.\ Phys.\ Lett.\ A \textbf{10}, 3077 (1995),
arXiv:hep-ph/9603369.

\bibitem{XingZhang2006}
Z.-Z.~Xing and H.~Zhang,
``On the Koide-like relations for the running masses of charged leptons,
neutrinos and quarks,''
Phys.\ Lett.\ B \textbf{635}, 107 (2006),
arXiv:hep-ph/0602134.

\bibitem{PDG2024}
Particle Data Group (S.~Navas \textit{et al.}),
``Review of Particle Physics,''
Phys.\ Rev.\ D \textbf{110}, 030001 (2024).

\bibitem{RiveroGsponer2005}
A.~Rivero and A.~Gsponer,
``The strange formula of Dr.\ Koide,''
arXiv:hep-ph/0505220 (2005).

\bibitem{ShahMir2007}
G.~N.~Shah and T.~A.~Mir,
``Pion and muon mass difference: a determining factor in elementary
particle mass distribution,''
Mod.\ Phys.\ Lett.\ A (2007),
arXiv:hep-ph/0702140.

\bibitem{LiMa2006}
N.~Li and B.-Q.~Ma,
``Energy scale independence of Koide's relation for quark and lepton masses,''
Phys.\ Rev.\ D \textbf{73}, 013009 (2006),
arXiv:hep-ph/0601031.

\bibitem{MasieroVeneziano1984}
A.~Masiero and G.~Veneziano,
``Split light composite supermultiplets,''
Nucl.\ Phys.\ B \textbf{249}, 593 (1985).

\bibitem{SanninoSchechter1997}
F.~Sannino and J.~Schechter,
``Chiral phase transition for SU(N) gauge theories via an effective
Lagrangian approach,''
Phys.\ Rev.\ D \textbf{60}, 056005 (1999),
arXiv:hep-ph/9903359.

\bibitem{Miyazawa1966}
H.~Miyazawa,
``Baryon number changing currents,''
Prog.\ Theor.\ Phys.\ \textbf{36}, 1266 (1966).

\end{thebibliography}

\end{document}
